% SECCIÓN 5: ALCANCE Y REQUISITOS
\section{Alcance y requisitos}

% SUB-SECCIÓN
\subsection{Alcance del proyecto}

	La intervención se limitará a 12 laboratorios de la EPIIS en la ciudad universitaria de Perayoc. De estos, 4 laboratorios recibirán intervención física y lógica para resolver problemas de lentitud de red, y 8 laboratorios recibirán exclusivamente intervención lógica para solucionar la caída del servicio DHCP en sus conexiones cableadas.

% SUB-SECCIÓN
\subsection{Perfil de carga y usuarios}

	\begin{itemize}
		\item \textbf{Capacidad por aula}: hasta 24 alumnos, operando al 100\% de capacidad en las horas pico (09:00 a 11:00 y 16:00 a 18:00). La ocupación en otros horarios varía entre el 50\% y el 70\%.
		\item \textbf{Usuarios principales}: estudiantes de diversas escuelas profesionales matriculados en el curso general \quotes{Pensamiento Computacional e Inteligencia Artificial}, demandando acceso a Google Classroom, Google Colab y Google Drive. Adicionalmente, estudiantes de Informática y Sistemas requiriendo acceso a la nube, desarrollo, bases de datos y servidores locales según sus mallas curriculares vigentes.
	\end{itemize}

% SUB-SECCIÓN
\subsection{Entregables y exclusiones}

	\textbf{Entregables.}

	\begin{itemize}
		\item Rediseño lógico implementado: VLANs por laboratorio y servicio DHCP optimizado (Capa 3) en los 8 laboratorios con fallas de asignación.
		\item Despliegue físico en los 4 laboratorios críticos: Access Points HP ARUBA AP-505 y switches Cisco Catalyst 9200L-24T-4X con cableado UTP Categoría 6.
		\item Pruebas de aceptación ejecutadas: RF-01 y RF-02 con aforo al 100\% y validación de acceso a servidores locales (RF-03).
		\item Documentación y manuales de la solución entregados.
	\end{itemize}

	\textbf{Exclusiones.}

	\begin{itemize}
		\item Los ambientes del pabellón distintos de los 12 laboratorios indicados.
		\item El enlace troncal de fibra óptica y la salida a internet de banda ancha, provistos por la UNSAAC.
		\item Las intervenciones sobre redes de datos móviles personales, que no forman parte de la red institucional evaluada.
		\item La ampliación de capacidad del enlace de la universidad; en caso de saturación se escalará a la oficina de TI central (ver riesgo R-02).
	\end{itemize}

% SUB-SECCIÓN
\subsection{EDT/WBS resumida}

	La estructura de desglose del trabajo sigue las fases del cronograma del apartado Viabilidad y cronograma:

	\begin{itemize}
		\item \textbf{1.0 Procesos de adquisición}: términos de referencia, cotizaciones, licitación, órdenes de compra y recepción de equipos.
		\item \textbf{2.0 Diagnóstico técnico}: monitoreo en horas pico, aislamiento del problema DHCP, inventario de puertos y prueba de estrés del enlace.
		\item \textbf{3.0 Rediseño lógico}: VLANs, servicio DHCP optimizado (Capa 3), ajuste de ruteo y PortFast.
		\item \textbf{4.0 Despliegue físico}: canalización e instalación de Access Points y switches en los 4 laboratorios críticos.
		\item \textbf{5.0 Pruebas y cierre}: pruebas de estrés al 100\%, validación de servidores locales y entrega de documentación y manuales.
	\end{itemize}

% SUB-SECCIÓN
\subsection{Requisitos funcionales}

	\begin{itemize}
		\item \textbf{RF-01 (Asignación IP)}: la red debe asignar una dirección IPv4 válida (DHCP) en menos de 5 segundos al conectar un equipo por cable en los 8 laboratorios afectados.
		\begin{itemize}
			\item \textit{Prueba}: conectar 24 equipos simultáneamente en un laboratorio y verificar la concesión IP y navegación exitosa.
		\end{itemize}
		\item \textbf{RF-02 (Conectividad inalámbrica)}: la red Wi-Fi en los 4 laboratorios críticos debe superar los anchos de banda actuales de 0.14 Mbps, permitiendo la ejecución fluida de entornos como Google Colab y carga de archivos a Drive.
		\begin{itemize}
			\item \textit{Prueba}: test de velocidad de subida/bajada y ejecución de un script de prueba en Google Colab con 24 usuarios concurrentes vía Wi-Fi.
		\end{itemize}
		\item \textbf{RF-03 (Acceso a servicios locales)}: la segmentación de red no debe impedir que los estudiantes de EPIIS accedan a los servidores locales requeridos para cursos de especialidad.
		\begin{itemize}
			\item \textit{Prueba}: ping y conexión exitosa por SSH/HTTP a los servidores de la facultad desde las VLANs de los laboratorios.
		\end{itemize}
	\end{itemize}

% SUB-SECCIÓN
\subsection{Requisitos no funcionales}

	\begin{itemize}
		\item \textbf{RNF-01 (Disponibilidad)}: las intervenciones físicas y los reinicios de red se ejecutan fuera del horario de clases pico (después de las 20:00 horas o fines de semana), para no interrumpir el servicio.
		\item \textbf{RNF-02 (Capacidad)}: la solución debe soportar 24 usuarios concurrentes por aula al 100\% de ocupación en horas pico.
		\item \textbf{RNF-03 (Compatibilidad)}: los equipos deben ser compatibles con los estándares ya utilizados en la escuela profesional.
		\item \textbf{RNF-04 (Segmentación segura)}: las VLANs no deben impedir el acceso de los estudiantes de EPIIS a los servidores locales de especialidad.
	\end{itemize}

% SUB-SECCIÓN
\subsection{Criterios de aceptación iniciales}

	Los requisitos funcionales se considerarán aceptados cuando se cumplan sus pruebas asociadas:

	\begin{itemize}
		\item \textbf{RF-01}: 24 equipos conectados simultáneamente por cable obtienen una dirección IPv4 por DHCP en menos de 5 segundos y navegan correctamente.
		\item \textbf{RF-02}: las pruebas de velocidad de subida y bajada superan los 0.14 Mbps actuales y un script de prueba en Google Colab se ejecuta con 24 usuarios concurrentes vía Wi-Fi.
		\item \textbf{RF-03}: se obtiene respuesta de ping y conexión exitosa por SSH/HTTP a los servidores de la facultad desde las VLANs de los laboratorios.
	\end{itemize}

	Los umbrales definitivos de velocidad y disponibilidad se fijarán después del diagnóstico de la Fase 1. \textit{[PENDIENTE: definir los umbrales objetivo de velocidad, latencia y disponibilidad.]}
